\section{The \texttt{Source\_div\_quasi} McXtrace Component}
Release: McXtrace 1.6

Quasi-stochastic X-ray source with Gaussian or uniform divergence

\subsection*{Identification}
\begin{itemize}
  \item \textbf{Author:} Mads Carlsen and Erik B Knudsen (erkn@fysik.dtu.dk)
  \item \textbf{Origin:} DTU Physics
  \item \textbf{Date:} Apr 21
\end{itemize}

\subsection*{Description}
A flat rectangular surface source with uniform or Gaussian divergence profile and focussing. If the parametere gauss is not set (the default) the divergence profile is flat in the range [-focus\_ax,focus\_ay]. If gauss is set, the focux\_ax,focus\_ay is considered the standard deviation of the gaussian profile. Currently focussing is only active for flat profile. The "focus window" is defined by focus\_xw,focus\_yh and dist. The spectral intensity profile is uniformly distributed in the energy interval defined by e0+-dE/2 or by wavelength lambda0+-dlambda/2

The phase space sapnned by the generated X-rays is sampled by means of Halton-sequences, instead of regular pseudo random numbers. This ensures that samples are evenly distributed within the phase space region of interest.

Example: Source\_div\_quasi(xwidth=0.1, yheight=0.1, focus\_aw=2, focus\_ah=2, E0=14, dE=2, gauss=0)

\subsection*{Input parameters}
Parameters in \textbf{boldface} are required; the others are optional.

\begin{longtable}{p{0.22\textwidth}p{0.12\textwidth}p{0.46\textwidth}p{0.14\textwidth}}
\toprule
\textbf{Name} & \textbf{Unit} & \textbf{Description} & \textbf{Default} \\
\midrule
\endhead
spectrum\_file & string & File from which to read the spectral intensity profile & "" \\
xwidth & m & Width of source & 0 \\
yheight & m & Height of source & 0 \\
focus\_xw & m & Width of sampling window & 0 \\
focus\_yh & m & Height of sampling window & 0 \\
dist & m & Downstream distance to place sampling target window & 0 \\
focus\_aw & rad & Std. dev. (Gaussian) or maximal (uniform) horz. width divergence. focus\_xw overrrides if it is more restrictive. & 0 \\
focus\_ah & rad & Std. dev. (Gaussian) or maximal (uniform) vert. height divergence. focus\_yh overrrides if it is more restrictive. & 0 \\
E0 & keV & Mean energy of X-rays. & 0 \\
dE & keV & Energy spread of X-rays. & 0 \\
lambda0 & \AA{} & Mean wavelength of X-rays (only relevant for E0=0) & 0 \\
dlambda & \AA{} & Wavelength half spread of X-rays. & 0 \\
flux & 1/(s*cm**2*st*energy unit) & Flux per energy unit, \AA{} or meV & 0 \\
gauss & 1 & Criterion: 0: uniform, 1: Gaussian distribution of energy/wavelength & 0 \\
gauss\_a & 1 & Criterion: 0: uniform, 1: Gaussian divergence distribution & 0 \\
randomphase & 0/1 & When=1, the X-ray phase is randomised & 1 \\
phase & rad & Set to finite value to define X-ray phase (0:2 pi) & 0 \\
verbose & 0/1 & Generate more output on the console. & 1 \\
\bottomrule
\end{longtable}

\subsection*{Links}
\begin{itemize}
  \item Component source code found in file \texttt{Source\_div\_quasi.comp}.
\end{itemize}
\IfFileExists{sources/Source_div_quasi_static.tex}{\input{sources/Source_div_quasi_static.tex}}{}